% ============================================================
%  Chapter 10 --- Write-Ahead Logging & Recovery
% ============================================================
\chapter{Write-Ahead Logging \& Recovery}
\label{ch:wal-recovery}

\section{The Durability Problem}

A DBMS must guarantee that committed transactions survive crashes.  But modern systems modify pages in the buffer pool (in memory) and only flush them to disk lazily.  If the system crashes between a commit and a flush, the committed data is lost.

\emph{Write-Ahead Logging} (WAL) solves this: before any page modification reaches disk, a log record describing the change is forced to persistent storage~\cite{mohan1992aries}.

\begin{conceptbox}[The WAL Rule]
\textbf{Before a modified page is written to disk, all log records pertaining to that page must be flushed to the log file.}  This guarantees that the log always contains enough information to redo or undo any in-flight change.
\end{conceptbox}

\section{Log Record Structure}

\begin{table}[H]
\centering
\caption{WAL log record fields.}
\begin{tabular}{l l p{8cm}}
\toprule
\textbf{Field} & \textbf{Type} & \textbf{Description} \\
\midrule
\code{LSN}           & \code{uint64\_t} & Log Sequence Number --- unique, monotonically increasing \\
\code{txnId}         & \code{uint32\_t} & Transaction that generated this record \\
\code{prevLSN}       & \code{uint64\_t} & Previous LSN in same transaction (for undo chain) \\
\code{type}          & enum             & \code{UPDATE}, \code{INSERT}, \code{DELETE}, \code{COMMIT}, \code{ABORT}, \code{CHECKPOINT}, \code{CLR} \\
\code{tableId}       & \code{uint32\_t} & Table affected \\
\code{pageId}        & \code{uint32\_t} & Page affected \\
\code{offset}        & \code{uint16\_t} & Byte offset within page \\
\code{beforeImage}   & \code{bytes}     & Old value (for undo) \\
\code{afterImage}    & \code{bytes}     & New value (for redo) \\
\bottomrule
\end{tabular}
\end{table}

\section{ARIES Recovery Algorithm}
\label{sec:aries}

ARIES~\cite{mohan1992aries} is the gold standard for WAL-based recovery.  It has three phases:

\begin{figure}[H]
\centering
\begin{tikzpicture}[node distance=0.4cm and 1cm]
  % Timeline
  \draw[-Stealth, thick] (0,0) -- (13,0) node[right] {Time};

  % Crash
  \draw[thick, shiftWarn] (10,0) -- (10,3.5);
  \node[font=\small\bfseries, shiftWarn] at (10,3.8) {CRASH};

  % Checkpoint
  \draw[thick, shiftAccent] (3,0) -- (3,2.5);
  \node[font=\small\bfseries, shiftAccent] at (3,2.8) {Checkpoint};

  % Phases
  \fill[shiftPrimary!10] (3,-0.5) rectangle (10,-1.5);
  \node[font=\small\bfseries] at (6.5,-1) {1. Analysis (scan forward)};

  \fill[shiftAccent!10] (3,-1.7) rectangle (10,-2.7);
  \node[font=\small\bfseries] at (6.5,-2.2) {2. Redo (replay all changes)};

  \fill[shiftWarn!10] (10,-2.9) rectangle (3,-3.9);
  \node[font=\small\bfseries] at (6.5,-3.4) {3. Undo (rollback losers)};

  \draw[-Stealth, shiftPrimary] (3.2,-1) -- (9.8,-1);
  \draw[-Stealth, shiftAccent] (3.2,-2.2) -- (9.8,-2.2);
  \draw[Stealth-, shiftWarn] (3.2,-3.4) -- (9.8,-3.4);
\end{tikzpicture}
\caption{The three phases of ARIES recovery.}
\label{fig:aries-phases}
\end{figure}

\subsection{Phase 1: Analysis}

\begin{algorithm}[H]
\caption{ARIES Analysis Phase}
\label{alg:aries-analysis}
\KwInput{Log starting from last checkpoint}
\KwOutput{Dirty Page Table (DPT), Active Transaction Table (ATT)}
Initialise DPT and ATT from checkpoint record\;
\ForEach{log record $r$ (forward scan)}{
  \If{$r$ is UPDATE/INSERT/DELETE}{
    Add $r.\text{pageId}$ to DPT if not present (with $r.\text{LSN}$ as \code{recLSN})\;
    Update ATT: mark $r.\text{txnId}$ as active\;
  }
  \If{$r$ is COMMIT}{
    Remove $r.\text{txnId}$ from ATT\;
  }
  \If{$r$ is ABORT}{
    Mark $r.\text{txnId}$ as ``undo'' in ATT\;
  }
}
\end{algorithm}

\subsection{Phase 2: Redo}

Replay \emph{all} logged changes from the earliest \code{recLSN} in the DPT, regardless of whether the transaction committed.  This restores the database to its exact pre-crash state.

\subsection{Phase 3: Undo}

Roll back all transactions that were active (not committed) at crash time.  Process undo records in reverse LSN order, writing \emph{Compensation Log Records} (CLRs) to ensure idempotency.

\section{Checkpointing}
\label{sec:checkpoint}

A checkpoint periodically records the state of the buffer pool and active transactions, limiting how far back ARIES must scan.

\begin{table}[H]
\centering
\caption{Checkpoint types.}
\begin{tabular}{l p{12cm}}
\toprule
\textbf{Type} & \textbf{Behaviour} \\
\midrule
Sharp (consistent) & Flush all dirty pages, halt all transactions during checkpoint.  Simple but causes a pause. \\
Fuzzy              & Write DPT and ATT to log without flushing pages.  No pause, but redo must start from earliest dirty page. \\
\bottomrule
\end{tabular}
\end{table}

\section{WAL in \shiftdb{} --- Design Path}

\shiftdb{} currently does not implement WAL.  Pages are flushed:
\begin{itemize}[nosep]
  \item On eviction from the buffer pool.
  \item On \code{BufferPool::flushAll()} at shutdown.
\end{itemize}

\begin{warningbox}[No WAL = No Crash Safety]
Without WAL, a crash after a \code{COMMIT} but before the dirty page is flushed will lose data.  Adding WAL would require:
\begin{enumerate}[nosep]
  \item A \code{LogManager} that appends log records to a sequential file.
  \item Modifying \code{BufferPool::markDirty} to write a log record first.
  \item Enforcing the WAL rule: flush log before flushing data page.
  \item Implementing a recovery procedure (ARIES three-phase) at startup.
\end{enumerate}
\end{warningbox}

\section{STEAL vs NO-STEAL, FORCE vs NO-FORCE}

\begin{table}[H]
\centering
\caption{Buffer management policy matrix.}
\begin{tabular}{l c c}
\toprule
 & \textbf{FORCE} & \textbf{NO-FORCE} \\
\midrule
\textbf{NO-STEAL} & No redo, no undo & Redo needed, no undo \\
\textbf{STEAL}     & No redo, undo needed & \textbf{Redo + Undo (ARIES)} \\
\bottomrule
\end{tabular}
\end{table}

\begin{itemize}[nosep]
  \item \textbf{STEAL}: the buffer pool can evict (``steal'') dirty pages from uncommitted transactions.  Requires \emph{undo} on crash.
  \item \textbf{NO-FORCE}: committed data does not have to be flushed to disk at commit time.  Requires \emph{redo} on crash.
  \item \textbf{ARIES} uses STEAL/NO-FORCE --- the most flexible and performant combination, but the most complex to implement.
\end{itemize}

\begin{interviewbox}[Interview Corner]
\textbf{Q:} What is the WAL principle and why is it necessary?\\[4pt]
\textbf{A:} WAL says ``log before data'' --- all changes must be logged to persistent storage before the corresponding data pages are written.  This ensures that on crash, the log contains enough information to redo committed changes and undo uncommitted ones.

\vspace{6pt}
\textbf{Q:} Describe the three phases of ARIES recovery.\\[4pt]
\textbf{A:} \emph{Analysis}: scan the log forward from the last checkpoint to identify dirty pages and active transactions.  \emph{Redo}: replay all logged changes to restore the database to its pre-crash state.  \emph{Undo}: roll back all uncommitted transactions in reverse order, writing CLRs.

\vspace{6pt}
\textbf{Q:} What is a Compensation Log Record (CLR)?\\[4pt]
\textbf{A:} A CLR is written during the undo phase to record the reversal of a logged change.  If the system crashes \emph{during recovery}, CLRs ensure that already-undone changes are not undone again.

\vspace{6pt}
\textbf{Q:} Explain STEAL/NO-FORCE and why ARIES uses this combination.\\[4pt]
\textbf{A:} STEAL allows dirty pages from uncommitted transactions to be flushed (good for memory management).  NO-FORCE avoids synchronous disk writes at commit (good for latency).  Together they provide maximum flexibility but require both redo and undo capabilities --- exactly what ARIES provides.
\end{interviewbox}
